% !TeX program = xelatex
% Main file. Edit the chapter files rather than this shared setup.
\documentclass[11pt,a4paper,oneside]{report}
\usepackage[T1]{fontenc}
\usepackage{lmodern}
\usepackage[margin=25mm,headheight=15pt]{geometry}
\usepackage{amsmath,amssymb,amsthm}
\usepackage{booktabs,enumitem,microtype}
\usepackage{xcolor,fancyhdr,needspace}
\usepackage{tikz}
\usetikzlibrary{arrows.meta,decorations.pathmorphing}
\usepackage{hyperref}
\definecolor{ink}{HTML}{16324F}
\definecolor{muted}{HTML}{52616B}
\definecolor{accent}{HTML}{A65332}
\hypersetup{colorlinks=true,linkcolor=ink,urlcolor=ink,
 pdftitle={Applied Mathematics 1 - Chapter Notes},bookmarksnumbered=true}
\setlength{\parindent}{0pt}
\setlength{\parskip}{5pt plus 1pt}
\setlength{\emergencystretch}{2em}
\setlist{itemsep=3pt,topsep=5pt}
\newtheorem{theorem}{Theorem}[chapter]
\newtheorem{proposition}[theorem]{Proposition}
\newtheorem{lemma}[theorem]{Lemma}
\theoremstyle{definition}
\newtheorem{definition}[theorem]{Definition}
\newtheorem{example}{Example}[chapter]
\theoremstyle{remark}
\newtheorem*{remark}{Remark}
\numberwithin{equation}{chapter}
\newcommand{\studyguide}[1]{\par\textbf{Method.} #1\par}
\pagestyle{fancy}
\fancyhf{}
\fancyhead[L]{\small Applied Mathematics 1}
\fancyhead[R]{\small\nouppercase{\leftmark}}
\fancyfoot[C]{\thepage}
\renewcommand{\chaptermark}[1]{\markboth{Chapter \thechapter}{} }
\setcounter{tocdepth}{1}
\title{Applied Mathematics 1\\[0.4em]\large}
\author{Roger D.}
\date{2026/27}
\begin{document}
\maketitle
\pagenumbering{roman}
\chapter*{How to use these notes}
\addcontentsline{toc}{chapter}{How to use these notes}
These notes are rewritten from the supplied MATH0008 typeset handouts and
handwritten lecture pages dated 5 October--6 November 2026. They are a
student study companion, not an official or verbatim set of lecture notes.
The original sources remain in \texttt{lecture\_resources/}.

The chapter order follows the handouts: discrete dynamical systems,
continuous models, then mechanics. Each model is presented through its
assumptions, equations, method of analysis and interpretation. Algebraic
steps omitted in the sources are filled in, and important qualifications
are stated explicitly. Additional explanations are labelled where they
go beyond the supplied derivation.

Source examples are identified where useful, but numbering in this edition
is automatic and need not match the originals.

Some typed material is ahead of the supplied handwritten pages: notably
parts of SIR analysis and the early mechanics examples. It is included as
handout material, not claimed to have been covered in a particular lecture.
Any optional or non-examinable label refers only to the supplied source;
check the current course for assessment scope.

Figures are newly drawn explanatory diagrams, not copies of source scans.
They are kept separately in \texttt{figures/} to keep the chapter source
readable. The README records the source inventory and principal corrections.

\textbf{Notation.} In a discrete model, $t=0,1,2,\ldots$ counts generations.
In a continuous model, $\dot x=dx/dt$ and $\ddot x=d^2x/dt^2$.
An equilibrium is denoted by a superscript $*$; a subscript $0$ denotes an
initial value. Populations and compartment sizes must be nonnegative.
Unless stated otherwise, model parameters described as rates are positive.
\tableofcontents
\clearpage
\pagenumbering{arabic}
% !TeX root = ../applied-mathematics1-companion.tex
\chapter{Discrete Dynamical Systems}\label{chap:applied-discrete}

\section{Building a discrete population model}
A discrete-time model describes a state at successive time steps:
\[
x_0,\ x_1,\ x_2,\ldots,\qquad x_{t+1}=f(x_t).
\]
It is appropriate, for example, for a population with distinct annual
breeding seasons. The initial condition $x_0$ must be specified.
Population values are treated as real expected sizes or large-population
approximations; the recurrence need not give integer outputs.
A physically meaningful population must satisfy $N_t\geq0$.

\subsection{Constant reproduction rate}
Suppose every individual produces on average $r\geq0$ offspring and the
old generation is replaced completely. Then
\[
N_{t+1}=rN_t,\qquad N_t=r^tN_0.
\]
Repeated substitution or induction proves the solution. If $N_0>0$:
\begin{itemize}
\item $0\leq r<1$: the population tends to zero;
\item $r=1$: the population stays constant;
\item $r>1$: it grows without bound.
\end{itemize}
For $r>0$, $r^t=e^{t\log r}$, hence the description \emph{exponential}
growth or decay. Time is still discrete even when the formula uses an
exponential function.

\subsection{Fixed harvesting}
If $h\geq0$ individuals are removed after reproduction, the algebraic model is
\[
N_{t+1}=rN_t-h.
\]
For $r\ne1$, iteration gives
\begin{align*}
N_t&=r^tN_0-h(1+r+\cdots+r^{t-1})\\
&=r^tN_0-h\frac{r^t-1}{r-1}\\
&=N^*+r^t(N_0-N^*),\qquad N^*=\frac h{r-1}.
\end{align*}
For $r=1$, $N_t=N_0-th$.

If $r>1$ and $h>0$, $N^*>0$ is a threshold:
\begin{itemize}
\item $N_0>N^*$ produces unbounded growth in this model;
\item $N_0=N^*$ produces a constant population;
\item $N_0<N^*$ eventually makes the algebraic formula nonpositive.
\end{itemize}
Negative population predictions indicate failure of the unmodified model.
The 5 October handwritten notes explicitly repair it as
\[
N_{t+1}=\max\{rN_t-h,0\}.
\]
Under this repaired model, a below-threshold population reaches zero in
finite time and remains extinct. The equality case $N_0=N^*$ is an
equilibrium, not extinction. For $h=0$, this reduces to reproduction alone.

\begin{example}[Checking a harvesting threshold]
As an added illustration, let $r=2$ and $h=10$. Then $N^*=10$.
Starting at $12$ gives $N_t=10+2^{t+1}$; starting at $10$ keeps it at $10$.
Starting at $8$ gives $8,6,2,0,0,\ldots$ in the clipped model.
The unclipped next value after $2$ would be $-6$, which is not a population.
\end{example}

\section{The logistic map and nondimensionalisation}
To represent reduced reproduction at high density, take
\[
N_{t+1}=rN_t\left(1-\frac{N_t}{K}\right),\qquad K>0.
\]
Set $x_t=N_t/K$. Dividing by $K$ gives the \emph{logistic map}
\begin{equation}\label{eq:applied-logistic-map}
x_{t+1}=rx_t(1-x_t)=f(x_t).
\end{equation}
Both $x_t$ and $r$ are dimensionless. Scaling removes $K$ from the
recurrence, leaving one parameter controlling the dynamics.

\begin{remark}[What $K$ means here]
In this discrete model, $K$ is the density scale where the next-generation
reproduction factor vanishes. The nonzero equilibrium is
$N^*=K(1-1/r)$, not $K$. This differs from the continuous logistic model,
whose carrying capacity is itself the positive equilibrium.
\end{remark}
For $0\leq x\leq1$, $0\leq x(1-x)\leq1/4$. Consequently, if
$0\leq r\leq4$, the interval $[0,1]$ is invariant: an initial value there
stays there. For $0<r<4$ and $0<x_0<1$, subsequent values remain in $(0,1)$.
For larger $r$, some initial values leave the biologically meaningful range.

\section{Fixed points and cobweb diagrams}
\begin{definition}[Fixed point]
A fixed point of $x_{t+1}=f(x_t)$ is a value $x^*$ with $f(x^*)=x^*$.
Starting exactly at $x^*$ gives a constant sequence. It is also called a
steady state or an equilibrium of the map.
\end{definition}
Fixed points are intersections of $y=f(x)$ with $y=x$.
For the logistic map, assuming $r>0$,
\[
x^*=rx^*(1-x^*)\quad\Longrightarrow\quad
x^*=0\quad\text{or}\quad x^*=1-\frac1r.
\]
The second point is positive only for $r>1$ and coincides with zero at $r=1$.

\studyguide{To construct a cobweb, plot $y=f(x)$ and $y=x$. From $x_0$
on the horizontal axis, go vertically to the curve: the height is $x_1$.
Go horizontally to the diagonal: the horizontal coordinate is now $x_1$.
Repeat vertically to the curve and horizontally to the diagonal.}
% !TeX root = ../applied-mathematics1-companion.tex
% Illustrative attracting affine map; the chapter explains general maps.
\begin{center}
\begin{tikzpicture}[scale=3.7,>=Stealth]
\draw[->] (0,0)--(1.2,0) node[right] {$x$};
\draw[->] (0,0)--(0,1.2) node[above] {$y$};
\draw[muted] (0,0)--(1.1,1.1) node[right] {$y=x$};
\draw[thick,ink] (0,.4)--(1.1,.95) node[right] {$y=f(x)=\tfrac12x+\tfrac25$};
\draw[thick,accent,->] (.1,0)--(.1,.45);
\draw[thick,accent,->] (.1,.45)--(.45,.45);
\draw[thick,accent,->] (.45,.45)--(.45,.625);
\draw[thick,accent,->] (.45,.625)--(.625,.625);
\draw[thick,accent] (.625,.625)--(.625,.7125)--(.7125,.7125)
 --(.7125,.75625)--(.75625,.75625)--(.75625,.778125)--(.778125,.778125);
\fill[ink] (.8,.8) circle (.012) node[above left] {$x^*=4/5$};
\node[below] at (.1,0) {$x_0$};
\end{tikzpicture}

{\small Cobweb steps: vertical to the map, horizontal to the diagonal.}
\end{center}


\section{Taylor approximation and local stability}
For a sufficiently differentiable function near $a$, a finite Taylor expansion
has the form
\[
f(a+h)=f(a)+f'(a)h+\frac12f''(a)h^2+O(h^3).
\]
Here $O(h^3)$ denotes an error bounded in magnitude by a constant times
$|h|^3$ for sufficiently small $h$, under appropriate smoothness assumptions.
In particular $f(a+h)=f(a)+f'(a)h+o(h)$ for a differentiable $f$.
The formal infinite Taylor series does not automatically equal every smooth
function; local approximation is all that is required for linearisation.

\begin{example}[Taylor expansion of $\log x$ near $2$]
Since $f'(x)=1/x$, $f''(x)=-1/x^2$ and $f'''(x)=2/x^3$,
\[
\log x=\log2+\frac{x-2}{2}-\frac{(x-2)^2}{8}
+\frac{(x-2)^3}{24}+O((x-2)^4).
\]
For a first-order approximation, keep only $\log2+(x-2)/2$.
\end{example}

\begin{definition}[Stability and attraction]
An equilibrium is \emph{stable} if arbitrarily small initial perturbations
can be kept arbitrarily small for all future times. It is \emph{locally
asymptotically stable} if it is stable and nearby states also converge to it.
An unstable equilibrium fails stability. These distinctions also apply to
continuous-time models.
\end{definition}
The sources often use ``stable'' to mean asymptotically stable. Here a
convergence conclusion is stated explicitly.

Write $x_t=x^*+y_t$. Since $f(x^*)=x^*$,
\[
y_{t+1}=f(x^*+y_t)-f(x^*)=f'(x^*)y_t+o(y_t).
\]
Discarding the smaller terms gives the \emph{linearised} recurrence
\[
y_{t+1}=\mu y_t,\qquad y_t=\mu^t y_0,\qquad \mu=f'(x^*).
\]
This explicit expression solves the linear approximation, not in general
the full nonlinear recurrence.

\begin{proposition}[Derivative test for a map]
For a continuously differentiable map near $x^*$:
\begin{itemize}
\item $|f'(x^*)|<1$ implies local asymptotic stability;
\item $|f'(x^*)|>1$ implies instability;
\item $|f'(x^*)|=1$ leaves the linear test inconclusive.
\end{itemize}
\end{proposition}
\begin{proof}[Explanation]
If $|f'(x^*)|<1$, continuity gives a neighbourhood with $|f'|\leq q<1$.
The mean value theorem yields $|f(x)-x^*|\leq q|x-x^*|$, so sufficiently
nearby iterates remain in the neighbourhood and converge geometrically.
If $|f'(x^*)|>1$, the same argument gives a local expansion factor greater
than one, so a nonzero perturbation eventually leaves a fixed small
neighbourhood. At equality, higher-order terms may decide the behaviour.
\end{proof}
For $0<\mu<1$, the linearised approach keeps the perturbation's sign.
For $-1<\mu<0$, it alternates sides while shrinking. At $\mu=0$,
higher-order terms determine the detailed approach, but the equilibrium
is still locally asymptotically stable.

\begin{example}[Logistic fixed-point stability]
Here $f'(x)=r(1-2x)$. At zero, $\mu=r$, so zero attracts for $0<r<1$
and is unstable for $r>1$. At $x^*=1-1/r$,
\[
\mu=r\left(1-2+\frac2r\right)=2-r.
\]
Thus the positive equilibrium attracts when $1<r<3$, and is unstable for
$r>3$. The approach is direct for $1<r<2$ and alternating for $2<r<3$.
The linear test alone does not settle $r=1$ or $r=3$.
\end{example}

\section{Parameters and bifurcations}
\begin{definition}[Bifurcation]
A bifurcation is a qualitative change in the dynamics as a parameter varies,
for example creation of equilibria or a change in their stability.
A bifurcation diagram plots equilibria or periodic states against a parameter;
solid and dashed lines often indicate attracting and unstable branches.
\end{definition}

\begin{example}[A parameter-dependent quadratic map]
For $x_{t+1}=x_t^2+x_t-a$, fixed points satisfy $(x^*)^2=a$.
For $a<0$ there are none; for $a=0$ there is one; for $a>0$ there are two:
$x^*=\pm\sqrt a$. Since $f'(x)=1+2x$,
\[
f'(\sqrt a)=1+2\sqrt a>1,\qquad
f'(-\sqrt a)=1-2\sqrt a.
\]
The positive branch is unstable for every $a>0$. The negative branch
attracts precisely in the range given by the strict derivative test,
$|1-2\sqrt a|<1$, that is $0<a<1$. It is unstable for $a>1$.
At $a=0$ and $a=1$, the multiplier equals $1$ or $-1$ respectively,
so that test is inconclusive at the parameter value itself.
These are the bifurcation values discussed in the sources.
\end{example}

\section{Period-two cycles and the route to chaos}
\begin{definition}[Period-two cycle]
Two distinct values $x_-,x_+$ form a period-two cycle if
$f(x_-)=x_+$ and $f(x_+)=x_-$. They satisfy
$f^{\circ2}(x)=x$, where $f^{\circ2}=f\circ f$ denotes composition, not
squaring the value $f(x)$.
\end{definition}
A fixed point of $f$ is also fixed by $f^{\circ2}$, so those old solutions
must be removed when looking for a genuine period-two cycle.

For the logistic map,
\[
f^{\circ2}(x)=r^2x(1-x)\bigl(1-rx(1-x)\bigr).
\]
Factoring $f^{\circ2}(x)-x=0$ gives the known roots $0,1-1/r$ and the
additional quadratic
\[
x^2-\left(1+\frac1r\right)x+\frac1r+\frac1{r^2}=0.
\]
Its two roots are
\begin{equation}\label{eq:applied-two-cycle}
x_\pm=\frac{r+1\pm\sqrt{r^2-2r-3}}{2r}.
\end{equation}
For $r>0$, they form a genuine distinct positive two-cycle when $r>3$.
At $r=3$ they merge with the fixed point; for the population interval,
we restrict attention to $r\leq4$.

\subsection{Stability of a cycle}
Let $g=f^{\circ2}$. By the chain rule, both cycle points have multiplier
\[
g'(x_-)=f'(x_+)f'(x_-)=g'(x_+).
\]
Using the sum and product of the two roots,
\begin{align*}
g'(x_\pm)
&=r^2(1-2x_-)(1-2x_+)\\
&=r^2\left[1-2\left(1+\frac1r\right)
+4\left(\frac1r+\frac1{r^2}\right)\right]\\
&=5-(r-1)^2.
\end{align*}
Hence the cycle attracts for
\[
|5-(r-1)^2|<1\quad\Longleftrightarrow\quad3<r<1+\sqrt6.
\]
A cycle can be attracting even though neither point is a fixed point of $f$:
we test the product of derivatives over a full cycle.

\subsection{Period doubling and chaotic behaviour}
At $r=3$, the fixed point loses attraction and the attracting two-cycle
appears. At $r=1+\sqrt6\approx3.44949$, the cycle multiplier reaches $-1$
and a four-cycle subsequently attracts. Further doublings give periods
$8,16,\ldots$, accumulating near $r\approx3.56995$.

Some larger parameters give bounded, aperiodic trajectories sensitive to
initial conditions: the source's introductory description of chaos.
Periodic windows also occur, including period three near
$1+\sqrt8\approx3.82843$; not every parameter beyond accumulation is chaotic.

% !TeX root = ../applied-mathematics1-companion.tex
\chapter{Continuous Mathematical Models}\label{chap:applied-continuous}

\section{From a verbal model to an ODE}
Continuous time describes a state through a rate of change:
\[
\dot x=\frac{dx}{dt}=F(x,t),\qquad x(0)=x_0.
\]
An equation is \emph{autonomous} when its right-hand side has no explicit
time dependence: $\dot x=f(x)$. Distinguish this from the discrete equation
$x_{t+1}=f(x_t)$: the same expression on the right can give very different
behaviour in the two settings.

\studyguide{Choose the state variables and units; state assumptions;
identify gains and losses; write rate of change = gains minus losses;
check units and physical constraints; impose initial conditions; then solve
or analyse equilibria and signs.}
For a scalar state of units $U$, the right side of an ODE must have units
$U/\text{time}$. This simple check often catches a missing factor.
The models here have smooth right sides, so standard local existence and
uniqueness results apply. A solution is followed only while it exists in
the relevant domain; not every nonlinear model has solutions for all time.

\subsection{Partial derivatives and first-order expansion}
For $f(x,y)$, partial differentiation varies one argument while fixing the
other. In symbols,
\[
f_x(x,y)=\lim_{h\to0}\frac{f(x+h,y)-f(x,y)}h,
\quad
f_y(x,y)=\lim_{h\to0}\frac{f(x,y+h)-f(x,y)}h.
\]
For instance, $f=x^2y+\cos y$ gives $f_x=2xy$ and $f_y=x^2-\sin y$.
For a differentiable function,
\[
f(a+X,b+Y)=f(a,b)+f_x(a,b)X+f_y(a,b)Y
+o\bigl(\sqrt{X^2+Y^2}\bigr).
\]
The derivatives evaluated at $(a,b)$ are constants in the linear approximation.

\section{Scalar population models}
\subsection{Exponential growth}
If growth rate is proportional to the population, $\dot P=kP$, $k>0$.
The solution is $P(t)=P_0e^{kt}$, including $P\equiv0$ when $P_0=0$.
Positive initial populations increase indefinitely. This can describe a
limited time range but ignores eventual resource constraints.

\subsection{Continuous logistic growth}
Competition within a population suggests
\begin{equation}\label{eq:applied-cont-logistic}
\dot P=kP\left(1-\frac PK\right),\qquad k,K>0.
\end{equation}
Here $k$ has units inverse time and $K$ has population units. The equilibria
are $0$ and $K$. For $0<P<K$, $\dot P>0$; for $P>K$, $\dot P<0$.
Thus the phase line points towards $K$ on both sides in the positive domain.

To solve, first treat the equilibria separately, since separation divides
by $P(K-P)$. For other initial values,
\begin{align*}
\frac{dP}{P(K-P)}&=\frac{k}{K}\,dt,\\
\frac1K\left(\frac1P+\frac1{K-P}\right)dP&=\frac{k}{K}\,dt,\\
\log\left|\frac P{K-P}\right|
-\log\left|\frac{P_0}{K-P_0}\right|&=kt.
\end{align*}
The sign of $P/(K-P)$ is preserved along a non-equilibrium trajectory.
After applying the initial condition and rearranging,
\begin{equation}\label{eq:applied-logistic-solution}
P(t)=\frac{KP_0}{P_0+(K-P_0)e^{-kt}}.
\end{equation}
This formula also gives both equilibrium solutions when $P_0=0$ or $K$.
For $P_0>0$, its denominator remains positive for $t\geq0$ and
$P(t)\to K$. If $P_0>K$, absolute values in the logarithmic derivation
are essential, even though the final formula has the same form.

\section{Phase lines and scalar stability}
For $\dot P=f(P)$, an equilibrium satisfies $f(P^*)=0$.
Write $P=P^*+p$ and expand:
\[
\dot p=f(P^*+p)=f'(P^*)p+o(p).
\]
The linearised equation $\dot p=\lambda p$ has solution
$p(t)=p_0e^{\lambda t}$, with $\lambda=f'(P^*)$.
For a continuously differentiable $f$ near $P^*$:
\begin{itemize}
\item $f'(P^*)<0$ implies local asymptotic stability;
\item $f'(P^*)>0$ implies instability;
\item $f'(P^*)=0$ leaves the linear test inconclusive.
\end{itemize}
The condition here is a \emph{negative derivative}, not derivative modulus
less than $1$. The latter belongs to discrete maps.

\begin{example}[A periodic array of equilibria: source Example 2.1]
For $\dot P=\sin P$, equilibria satisfy $P^*=n\pi$, $n\in\mathbb Z$.
Since $f'(P^*)=\cos(n\pi)=(-1)^n$, odd multiples attract and even multiples
are unstable. Between neighbouring even multiples $2j\pi$ and
$2(j+1)\pi$, every non-equilibrium solution tends to $(2j+1)\pi$.
If $P$ denotes a population, restrict the mathematical phase line to
nonnegative values.
\end{example}

\begin{example}[Additional handwritten example, 15 October]
Consider $\dot x=f(x)=e^{-x}\sin x$ for initial values $x_0\in[0,6]$.
The equilibria in this interval are $0$ and $\pi$.
Since $f'(x)=e^{-x}(\cos x-\sin x)$,
\[
f'(0)=1>0,\qquad f'(\pi)=-e^{-\pi}<0.
\]
Thus $0$ is unstable, while $\pi$ attracts. For $0<x<\pi$ the state moves
right; for $\pi<x\leq6<2\pi$ it moves left. Every $x_0\in(0,6]$ tends to
$\pi$, while $x_0=0$ stays fixed.

To interpret the changing speed, note $|\dot x|=|f(x)|$.
The positive velocity has its maximum at $x=\pi/4$; the negative velocity
has its minimum at $x=5\pi/4$. A right-moving state below $\pi/4$ first
speeds up and then slows down. A left-moving state above $5\pi/4$ first
speeds up as it approaches $5\pi/4$, then slows towards $\pi$.
A negative velocity and a decreasing speed are different statements.
\end{example}

\section{A mixing model}
A tank initially contains $100$ litres of pure water. Brine with concentration
$5$ grams per litre enters at $2$ litres per minute; an equal outflow keeps
the volume constant. Assume perfect instantaneous mixing.
Let $M(t)$ be the salt mass in grams and $t$ be measured in minutes.
\[
\text{inflow salt rate}=5\times2=10\ \mathrm{g/min},\qquad
\text{outflow salt rate}=\frac{M}{100}\times2=\frac{M}{50}\ \mathrm{g/min}.
\]
Therefore $\dot M=10-M/50$, with $M(0)=0$. An integrating factor gives
\begin{align*}
\frac d{dt}\left(Me^{t/50}\right)&=10e^{t/50},\\
M(t)&=500+Ce^{-t/50},\\
M(0)=0&\quad\Longrightarrow\quad C=-500.
\end{align*}
The mass and concentration are
\[
M(t)=500(1-e^{-t/50})\ \mathrm g,\qquad
c(t)=\frac{M(t)}{100}=5(1-e^{-t/50})\ \mathrm{g/L}.
\]
As $t\to\infty$, the tank concentration approaches the incoming concentration.
If inflow and outflow volumes differed, the volume would vary and the
outflow concentration would be $M(t)/V(t)$ instead.

\section{Coupled systems, phase planes and nullclines}
A two-species model has the form
\[
\dot x=f(x,y),\qquad\dot y=g(x,y).
\]
Terms proportional to $xy$ represent encounters under a well-mixed
assumption. Their signs encode the interaction. For positive constants:
\begin{center}
\begin{tabular}{lll}
\toprule
Interaction&$\dot x$&$\dot y$\\\midrule
Predation ($x$ prey, $y$ predator)&$ax-bxy$&$-cy+dxy$\\
Competition&$ax-bxy$&$cy-dxy$\\
Mutual benefit&$ax+bxy$&$cy+dxy$\\\bottomrule
\end{tabular}
\end{center}
These are simplified models; sign information alone does not capture all
resource limitations or guarantee bounded populations.

\begin{definition}[Phase portrait and nullclines]
A trajectory is the curve traced by $(x(t),y(t))$ in the phase plane,
with arrows showing increasing time. The vector field is $(f(x,y),g(x,y))$.
The $x$-nullclines satisfy $f(x,y)=0$, where nonzero arrows are vertical.
The $y$-nullclines satisfy $g(x,y)=0$, where nonzero arrows are horizontal.
Equilibria satisfy both equations simultaneously.
\end{definition}
Under uniqueness, two distinct trajectories cannot cross at a non-equilibrium
point: the same state cannot have two different future solutions.
A nullcline is generally not a trajectory, and a point on just one
nullcline need not be an equilibrium.

\begin{example}[Reading the source phase portrait: source Example 2.3]
The supplied sketch has rightward flow on the $x$-axis and downward flow
on the $y$-axis, with an unstable origin, an unstable interior equilibrium,
and a closed orbit attracting nearby trajectories. It illustrates an
attracting \emph{limit cycle}: an isolated periodic orbit.
One cannot infer this entire portrait from interaction signs alone.
In particular, the closed Lotka--Volterra orbits studied below are not
attracting limit cycles, because they form a family of neighbouring cycles.
\end{example}

\studyguide{For a phase portrait: solve $f=g=0$; draw both types of nullcline;
check signs of $f$ and $g$ in each region; combine them into flow arrows;
then draw trajectories consistent with those directions and uniqueness.}

\begin{example}[Competing species: source Section 2.4.2]\label{ex:applied-competition}
Consider
\[
\dot x=x(5-x-y),\qquad\dot y=y(8-x-2y).
\]
The $x$-nullclines are $x=0$ and $x+y=5$; the $y$-nullclines are $y=0$
and $x+2y=8$. Their intersections give
\[
(0,0),\qquad(5,0),\qquad(0,4),\qquad(2,3).
\]
For $x,y>0$, $\dot x$ has the sign of $5-x-y$, and $\dot y$ has the sign
of $8-x-2y$. At $(1,1)$, both are positive, so arrows point up and right.
Above both oblique nullclines, both are negative, so arrows point down and left.
The axes are invariant because the corresponding population derivative vanishes.
\end{example}
% !TeX root = ../applied-mathematics1-companion.tex
\begin{center}
\begin{tikzpicture}[x=.8cm,y=.8cm,>=Stealth]
\draw[->] (0,0)--(8.8,0) node[right] {$x$};
\draw[->] (0,0)--(0,5.8) node[above] {$y$};
\foreach \x in {.5,1.5,2.5,3.5,4.5,5.5,6.5,7.5}
 \foreach \y in {.5,1.5,2.5,3.5,4.5} {
 \pgfmathsetmacro{\dx}{\x*(5-\x-\y)}
 \pgfmathsetmacro{\dy}{\y*(8-\x-2*\y)}
 \pgfmathsetmacro{\norm}{max(sqrt(\dx*\dx+\dy*\dy),.01)}
 \draw[->,muted!60] (\x,\y)--++({.24*\dx/\norm},{.24*\dy/\norm});}
\draw[thick,ink] (0,5)--(5,0);
\draw[thick,accent,dashed] (0,4)--(8,0);
\node[ink,above right] at (.2,4.8) {$x+y=5$};
\node[accent,above right] at (5,1.5) {$x+2y=8$};
\foreach \p in {(0,0),(5,0),(0,4),(2,3)} \fill \p circle (2pt);
\node[below left] at (0,0) {$0$};
\node[below] at (5,0) {$5$};
\node[left] at (0,4) {$4$};
\node[above right,fill=white,inner sep=1pt] at (2,3) {$(2,3)$};
\end{tikzpicture}

{\small Competition model: oblique nullclines and normalised flow arrows.}
\end{center}


\section{Linearisation of two-dimensional systems}
At an equilibrium $(x^*,y^*)$, write $x=x^*+X$ and $y=y^*+Y$.
First-order expansion gives the linearised system
\[
\begin{pmatrix}\dot X\\\dot Y\end{pmatrix}
=J\begin{pmatrix}X\\Y\end{pmatrix},\qquad
J=\begin{pmatrix}A&B\\C&D\end{pmatrix}
=\left.\begin{pmatrix}f_x&f_y\\g_x&g_y\end{pmatrix}\right|_{(x^*,y^*)}.
\]
The matrix $J$ is the \emph{Jacobian at the equilibrium}. For the original
nonlinear system there are also higher-order remainder terms.

Try a mode $(X,Y)^{\mathsf T}=\mathbf v e^{\lambda t}$, with
$\mathbf v\ne\mathbf0$. Substitution gives $J\mathbf v=\lambda\mathbf v$.
Thus $\lambda$ is an eigenvalue and must satisfy
\begin{equation}\label{eq:applied-characteristic}
\det(\lambda I-J)=(\lambda-A)(\lambda-D)-BC=0.
\end{equation}
This determinant approach works even when $B=0$ or $C=0$; dividing by a
matrix entry could lose those cases.

For distinct eigenvalues with independent eigenvectors, the general linear
solution is
\[
\begin{pmatrix}X\\Y\end{pmatrix}
=c_1\mathbf v_1e^{\lambda_1t}+c_2\mathbf v_2e^{\lambda_2t}.
\]
There are two independent constants, not four unrelated coefficients in
$X$ and $Y$. Repeated eigenvalues may require a term involving $te^{\lambda t}$.
Complex conjugate modes combine into real sine and cosine terms.
Superposition applies to the homogeneous linearised system, not generally
to the original nonlinear model.

\begin{proposition}[Local linear stability]
For a continuously differentiable nonlinear vector field, if all eigenvalues
of its equilibrium Jacobian have negative real part, the equilibrium is
locally asymptotically stable. If at least one has positive real part,
it is unstable. If some real parts are zero and none are positive,
linearisation alone generally does not decide nonlinear stability.
\end{proposition}
These standard linearisation results justify using the modes as a local
stability test. The linear solution is not an exact formula for the full
nonlinear trajectory, even when it predicts stability correctly.

\subsection{Trace and determinant: a supplementary shortcut}
Let $\tau=A+D$ and $\Delta=AD-BC$. The characteristic equation is
$\lambda^2-\tau\lambda+\Delta=0$.
\begin{itemize}
\item $\Delta<0$: one positive and one negative eigenvalue, an unstable saddle.
\item $\Delta>0$, $\tau<0$: negative real parts, local attraction.
\item $\Delta>0$, $\tau>0$: positive real parts, instability.
\item $\Delta>0$, $\tau=0$: purely imaginary eigenvalues; no nonlinear
stability conclusion from the linear test alone.
\end{itemize}
For $\Delta>0$, the sign of $\tau^2-4\Delta$ distinguishes two real modes
from complex conjugate modes. The cases $\Delta=0$ require further analysis.

\subsection{Return to the competition example}
At $(2,3)$, direct substitution gives
\begin{align*}
\dot X&=-2X-2Y-X^2-XY,\\
\dot Y&=-3X-6Y-XY-2Y^2.
\end{align*}
Dropping the quadratic terms gives
\[
J=\begin{pmatrix}-2&-2\\-3&-6\end{pmatrix},\qquad
\lambda^2+8\lambda+6=0,\qquad\lambda=-4\pm\sqrt{10}.
\]
Both eigenvalues are negative, so coexistence is locally asymptotically stable.
For reference, the general Jacobian is
\[
J(x,y)=\begin{pmatrix}5-2x-y&-x\\-y&8-x-4y\end{pmatrix}.
\]
The other equilibria have eigenvalues $5,8$ at $(0,0)$; $-5,3$ at $(5,0)$;
and $1,-8$ at $(0,4)$. Thus the origin is repelling and the two nonzero
axis equilibria are saddles. These extra classifications complete the
local picture without asserting global attraction for every initial state.

\section{The Lotka--Volterra predator--prey model}
\textbf{Convention for this section: $x$ is predator population, $y$ is prey.}
This reverses the labels in the earlier illustrative interaction table.
With $a,b,c,d>0$,
\begin{equation}\label{eq:applied-lv}
\dot x=x(-a+by),\qquad\dot y=y(c-dx).
\end{equation}
Predators decline without prey, prey grow without predators, and encounters
benefit predators while reducing prey. Positive populations stay positive:
for example $x(t)=x_0\exp\!\left(\int_0^t[-a+by(s)]\,ds\right)$.
The axes are invariant and are analysed separately from the positive interior.

The nullclines are $x=0$, $y=a/b$ for $\dot x=0$, and $y=0$, $x=c/d$
for $\dot y=0$. Equilibria are $(0,0)$ and $(c/d,a/b)$.
Above $y=a/b$ predators grow; to the right of $x=c/d$ prey decrease.
For this predator-on-horizontal-axis convention, the positive orbits
run clockwise around coexistence.

\subsection{Harvesting}
Per-capita harvesting rates $h_1,h_2\geq0$ give
\[
\dot x=x(-a-h_1+by),\qquad\dot y=y(c-h_2-dx).
\]
If $c>h_2$, the positive equilibrium shifts to
\[
\left(\frac{c-h_2}{d},\frac{a+h_1}{b}\right).
\]
Thus its predator value falls and prey value rises relative to the original
coexistence equilibrium. These are equilibrium or cycle-average comparisons,
not a claim that every instantaneous prey population rises. If $c\leq h_2$,
there is no positive coexistence equilibrium of this form.

\subsection{Local analysis}
At the origin, the linearised equations are $\dot X=-aX$, $\dot Y=cY$,
so the origin is a saddle.
At coexistence,
\[
J=\begin{pmatrix}0&bc/d\\-ad/b&0\end{pmatrix},\qquad
\lambda^2=-ac,\qquad\lambda=\pm i\sqrt{ac}.
\]
Let $\omega=\sqrt{ac}$ and $B=bc/d$. The real linear solution can be written
\begin{align*}
X(t)&=A_1\cos(\omega t)+A_2\sin(\omega t),\\
Y(t)&=\frac\omega B\bigl[-A_1\sin(\omega t)+A_2\cos(\omega t)\bigr].
\end{align*}
It oscillates with period $2\pi/\sqrt{ac}$, using two independent constants.
This is the small-amplitude linearised period, not the exact period of
every nonlinear orbit. Purely imaginary eigenvalues alone do not prove
nonlinear stability; we next use a conserved quantity.

\subsection{A first integral and closed trajectories}
For $x,y>0$ and away from a vertical trajectory tangent, take the ratio:
\[
\frac{dy}{dx}=\frac{y(c-dx)}{x(-a+by)}.
\]
Separation yields
\[
\int\left(-\frac ay+b\right)dy
=\int\left(\frac cx-d\right)dx,
\]
so a trajectory satisfies
\begin{equation}\label{eq:applied-lv-integral}
c\log x-dx+a\log y-by=\text{constant}.
\end{equation}
Logs require \emph{strictly positive} populations, not merely nonnegative ones.
Direct differentiation verifies the conserved expression along the whole
positive trajectory, including points where the ratio $dy/dx$ was unavailable.
Exponentiating gives
\[
F(x)G(y)=H>0,\qquad F(x)=x^ce^{-dx},\quad G(y)=y^ae^{-by}.
\]
Each factor tends to zero at both ends of its positive domain. Also
\[
F'(x)=x^{c-1}e^{-dx}(c-dx),\qquad
G'(y)=y^{a-1}e^{-by}(a-by).
\]
Thus $F$ has a unique maximum $M_1=(c/d)^ce^{-c}$ at $c/d$;
$G$ has a unique maximum $M_2=(a/b)^ae^{-a}$ at $a/b$.
At $H=M_1M_2$, the level set is just the equilibrium. For
$0<H<M_1M_2$, horizontal and vertical lines through the equilibrium each
meet the level set twice, and the trajectory is a closed curve surrounding it.

\textit{Supplementary justification.} Equivalently,
\[
V(x,y)=dx-c\log x+by-a\log y
\]
is strictly convex, tends to $+\infty$ at the boundary and at infinity,
and has its unique minimum at coexistence. Its non-minimal level sets
are compact smooth closed curves with no equilibrium on them.
The nonvanishing vector field follows each such curve periodically.
It also follows that coexistence is stable in the non-attracting sense:
nearby level curves remain nearby, but no other one converges to the centre.
It is therefore a nonlinear \emph{centre}, not an asymptotically stable state
and not an attracting limit cycle.
% !TeX root = ../applied-mathematics1-companion.tex
% Newly integrated illustrative trajectories for a=b=c=d=1.
\begin{center}
\begin{tikzpicture}[x=1.8cm,y=1.8cm,>=Stealth]
\draw[->] (0,0)--(3.1,0) node[right] {$x$ (predators)};
\draw[->] (0,0)--(0,3.1) node[above] {$y$ (prey)};
\draw[muted,dashed] (1,0)--(1,2.9);
\draw[muted,dashed] (0,1)--(2.9,1);
\draw[ink,thick] plot[smooth] coordinates {(1.3000,1.0000) (1.3000,0.9988) (1.2995,0.9845) (1.2981,0.9705) (1.2958,0.9568) (1.2927,0.9433) (1.2888,0.9303) (1.2841,0.9176) (1.2786,0.9053) (1.2725,0.8934) (1.2656,0.8819) (1.2581,0.8709) (1.2500,0.8603) (1.2414,0.8502) (1.2322,0.8406) (1.2225,0.8315) (1.2124,0.8229) (1.2019,0.8147) (1.1911,0.8071) (1.1799,0.7999) (1.1684,0.7932) (1.1567,0.7871) (1.1448,0.7814) (1.1327,0.7762) (1.1205,0.7715) (1.1081,0.7673) (1.0957,0.7635) (1.0833,0.7603) (1.0708,0.7575) (1.0583,0.7551) (1.0459,0.7532) (1.0336,0.7518) (1.0213,0.7508) (1.0092,0.7503) (0.9971,0.7502) (0.9852,0.7505) (0.9735,0.7512) (0.9620,0.7524) (0.9507,0.7540) (0.9395,0.7559) (0.9287,0.7583) (0.9180,0.7611) (0.9076,0.7643) (0.8975,0.7679) (0.8876,0.7719) (0.8780,0.7762) (0.8688,0.7810) (0.8598,0.7861) (0.8511,0.7916) (0.8427,0.7974) (0.8347,0.8036) (0.8270,0.8101) (0.8196,0.8170) (0.8126,0.8243) (0.8059,0.8319) (0.7996,0.8398) (0.7936,0.8480) (0.7880,0.8566) (0.7828,0.8655) (0.7779,0.8746) (0.7734,0.8841) (0.7693,0.8939) (0.7656,0.9039) (0.7623,0.9142) (0.7593,0.9247) (0.7568,0.9356) (0.7546,0.9466) (0.7529,0.9578) (0.7516,0.9693) (0.7507,0.9810) (0.7502,0.9928) (0.7502,1.0048) (0.7506,1.0169) (0.7514,1.0291) (0.7527,1.0414) (0.7544,1.0538) (0.7566,1.0663) (0.7592,1.0787) (0.7623,1.0912) (0.7659,1.1036) (0.7699,1.1160) (0.7744,1.1283) (0.7795,1.1404) (0.7850,1.1524) (0.7909,1.1642) (0.7974,1.1757) (0.8044,1.1870) (0.8119,1.1980) (0.8198,1.2087) (0.8283,1.2189) (0.8372,1.2287) (0.8467,1.2381) (0.8566,1.2469) (0.8670,1.2552) (0.8778,1.2630) (0.8891,1.2700) (0.9009,1.2765) (0.9130,1.2822) (0.9256,1.2872) (0.9385,1.2914) (0.9518,1.2948) (0.9655,1.2973) (0.9794,1.2991) (0.9936,1.2999) (1.0080,1.2999) (1.0226,1.2989) (1.0373,1.2970) (1.0521,1.2943) (1.0670,1.2906) (1.0819,1.2860) (1.0967,1.2805) (1.1114,1.2741) (1.1259,1.2668) (1.1402,1.2588) (1.1542,1.2499) (1.1679,1.2403) (1.1812,1.2299) (1.1939,1.2189) (1.2062,1.2073) (1.2179,1.1950) (1.2290,1.1823) (1.2394,1.1691) (1.2491,1.1554) (1.2580,1.1414) (1.2662,1.1272) (1.2735,1.1127) (1.2799,1.0980) (1.2855,1.0832) (1.2902,1.0683) (1.2940,1.0534) (1.2968,1.0386) (1.2988,1.0238) (1.2998,1.0092) (1.3000,1.0000)};
\draw[->,ink,thick] (1.30,1)--(1.30,.91);
\draw[ink,thick] plot[smooth] coordinates {(1.7000,1.0000) (1.7000,0.9972) (1.6984,0.9643) (1.6942,0.9326) (1.6875,0.9021) (1.6784,0.8730) (1.6671,0.8453) (1.6537,0.8189) (1.6384,0.7939) (1.6213,0.7703) (1.6027,0.7480) (1.5826,0.7270) (1.5613,0.7073) (1.5388,0.6889) (1.5153,0.6716) (1.4911,0.6556) (1.4661,0.6407) (1.4405,0.6269) (1.4145,0.6142) (1.3882,0.6025) (1.3616,0.5917) (1.3348,0.5819) (1.3080,0.5730) (1.2813,0.5650) (1.2546,0.5577) (1.2280,0.5513) (1.2017,0.5457) (1.1756,0.5407) (1.1499,0.5365) (1.1245,0.5330) (1.0995,0.5302) (1.0749,0.5280) (1.0508,0.5264) (1.0271,0.5254) (1.0040,0.5250) (0.9813,0.5252) (0.9592,0.5259) (0.9377,0.5272) (0.9167,0.5291) (0.8962,0.5315) (0.8764,0.5344) (0.8571,0.5378) (0.8383,0.5417) (0.8202,0.5462) (0.8026,0.5512) (0.7856,0.5567) (0.7692,0.5626) (0.7533,0.5691) (0.7380,0.5761) (0.7233,0.5836) (0.7091,0.5916) (0.6955,0.6001) (0.6824,0.6092) (0.6699,0.6187) (0.6579,0.6288) (0.6464,0.6394) (0.6355,0.6505) (0.6251,0.6621) (0.6152,0.6743) (0.6059,0.6870) (0.5970,0.7003) (0.5887,0.7141) (0.5809,0.7285) (0.5735,0.7434) (0.5667,0.7589) (0.5604,0.7750) (0.5546,0.7916) (0.5493,0.8089) (0.5446,0.8266) (0.5403,0.8450) (0.5365,0.8639) (0.5333,0.8834) (0.5305,0.9035) (0.5284,0.9242) (0.5267,0.9454) (0.5256,0.9671) (0.5250,0.9894) (0.5251,1.0122) (0.5257,1.0355) (0.5269,1.0594) (0.5287,1.0836) (0.5311,1.1084) (0.5342,1.1335) (0.5380,1.1591) (0.5424,1.1849) (0.5476,1.2111) (0.5535,1.2375) (0.5602,1.2641) (0.5678,1.2908) (0.5761,1.3177) (0.5853,1.3444) (0.5955,1.3711) (0.6066,1.3977) (0.6186,1.4239) (0.6317,1.4498) (0.6459,1.4751) (0.6612,1.4999) (0.6777,1.5238) (0.6953,1.5470) (0.7142,1.5690) (0.7343,1.5900) (0.7558,1.6095) (0.7786,1.6276) (0.8027,1.6441) (0.8282,1.6587) (0.8551,1.6714) (0.8833,1.6819) (0.9129,1.6902) (0.9438,1.6960) (0.9759,1.6993) (1.0093,1.6999) (1.0437,1.6977) (1.0792,1.6927) (1.1154,1.6849) (1.1524,1.6741) (1.1900,1.6604) (1.2278,1.6438) (1.2658,1.6245) (1.3036,1.6024) (1.3411,1.5778) (1.3779,1.5508) (1.4138,1.5216) (1.4486,1.4904) (1.4820,1.4575) (1.5136,1.4231) (1.5434,1.3874) (1.5710,1.3508) (1.5962,1.3135) (1.6189,1.2757) (1.6390,1.2378) (1.6563,1.1999) (1.6708,1.1622) (1.6823,1.1251) (1.6910,1.0886) (1.6967,1.0529) (1.6996,1.0182) (1.7000,1.0000)};
\draw[->,ink,thick] (1.70,1)--(1.70,.91);
\draw[ink,thick] plot[smooth] coordinates {(2.3000,1.0000) (2.3000,0.9948) (2.2961,0.9347) (2.2857,0.8785) (2.2696,0.8262) (2.2481,0.7778) (2.2218,0.7330) (2.1913,0.6917) (2.1571,0.6538) (2.1198,0.6190) (2.0798,0.5872) (2.0375,0.5581) (1.9935,0.5316) (1.9480,0.5074) (1.9014,0.4853) (1.8541,0.4653) (1.8064,0.4471) (1.7583,0.4306) (1.7103,0.4157) (1.6625,0.4023) (1.6150,0.3901) (1.5680,0.3792) (1.5216,0.3694) (1.4759,0.3607) (1.4310,0.3529) (1.3870,0.3460) (1.3439,0.3400) (1.3019,0.3348) (1.2608,0.3303) (1.2208,0.3265) (1.1819,0.3234) (1.1441,0.3209) (1.1073,0.3189) (1.0717,0.3176) (1.0371,0.3167) (1.0036,0.3164) (0.9712,0.3166) (0.9399,0.3173) (0.9096,0.3184) (0.8804,0.3201) (0.8522,0.3221) (0.8249,0.3246) (0.7987,0.3276) (0.7734,0.3310) (0.7490,0.3348) (0.7255,0.3390) (0.7029,0.3437) (0.6812,0.3488) (0.6603,0.3544) (0.6403,0.3604) (0.6210,0.3668) (0.6025,0.3737) (0.5848,0.3811) (0.5678,0.3889) (0.5515,0.3972) (0.5359,0.4060) (0.5209,0.4153) (0.5066,0.4251) (0.4930,0.4355) (0.4799,0.4463) (0.4674,0.4578) (0.4556,0.4698) (0.4442,0.4823) (0.4335,0.4955) (0.4232,0.5093) (0.4135,0.5237) (0.4043,0.5388) (0.3956,0.5545) (0.3874,0.5709) (0.3797,0.5881) (0.3724,0.6060) (0.3656,0.6246) (0.3592,0.6440) (0.3533,0.6642) (0.3478,0.6853) (0.3428,0.7071) (0.3382,0.7299) (0.3340,0.7535) (0.3303,0.7781) (0.3270,0.8036) (0.3241,0.8300) (0.3217,0.8574) (0.3197,0.8858) (0.3182,0.9153) (0.3171,0.9458) (0.3165,0.9773) (0.3164,1.0099) (0.3169,1.0436) (0.3178,1.0783) (0.3193,1.1142) (0.3213,1.1511) (0.3239,1.1892) (0.3272,1.2283) (0.3311,1.2685) (0.3357,1.3098) (0.3411,1.3520) (0.3473,1.3953) (0.3543,1.4394) (0.3622,1.4845) (0.3712,1.5303) (0.3812,1.5768) (0.3923,1.6239) (0.4047,1.6715) (0.4184,1.7194) (0.4336,1.7674) (0.4504,1.8154) (0.4689,1.8631) (0.4893,1.9103) (0.5118,1.9567) (0.5364,2.0019) (0.5634,2.0457) (0.5930,2.0875) (0.6254,2.1271) (0.6608,2.1639) (0.6993,2.1974) (0.7412,2.2271) (0.7867,2.2525) (0.8358,2.2731) (0.8889,2.2882) (0.9458,2.2973) (1.0066,2.3000) (1.0714,2.2957) (1.1398,2.2841) (1.2118,2.2650) (1.2868,2.2380) (1.3645,2.2033) (1.4443,2.1610) (1.5252,2.1113) (1.6067,2.0547) (1.6876,1.9919) (1.7670,1.9235) (1.8439,1.8505) (1.9173,1.7739) (1.9860,1.6947) (2.0494,1.6139) (2.1066,1.5324) (2.1569,1.4514) (2.1999,1.3715) (2.2353,1.2936) (2.2629,1.2183) (2.2827,1.1460) (2.2950,1.0773) (2.2999,1.0122) (2.3000,1.0000)};
\draw[->,ink,thick] (2.30,1)--(2.30,.91);
\fill (1,1) circle (1.5pt);
\node[below] at (1,0) {$c/d=1$};
\node[left] at (0,1) {$a/b=1$};
\end{tikzpicture}

{\small Closed orbits around the centre $(1,1)$, for $a=b=c=d=1$.}
\end{center}


\subsection{Average populations over a cycle}
For a periodic positive orbit of period $T$,
\[
\frac{\dot y}y=c-dx\quad\Longrightarrow\quad
\frac1T\int_0^T x(t)\,dt
=\frac cd-\frac{\log y(T)-\log y(0)}{dT}=\frac cd.
\]
Similarly, integrating $\dot x/x=-a+by$ gives
\[
\overline x=\frac cd,\qquad\overline y=\frac ab.
\]
The time averages equal the equilibrium populations although the instantaneous
values oscillate. For the harvested model with positive coexistence, the
same calculation uses $a+h_1$ and $c-h_2$.

\section{The SIR epidemic model}
This is a mathematical compartment model under idealised assumptions,
not a fitted description of a particular current epidemic.
Let $S$ denote susceptible individuals, $I$ infectious individuals, and
$R$ removed individuals (recovered, isolated or dead).
Assume a closed, well-mixed population, no demographic turnover or return
to susceptibility, and transitions $S\to I\to R$.
The model is
\begin{equation}\label{eq:applied-sir}
\dot S=-\alpha SI,\qquad
\dot I=\alpha SI-\beta I,\qquad
\dot R=\beta I,\qquad\alpha,\beta>0.
\end{equation}
With population counts as variables, $\alpha$ has units
$(\text{population}\times\text{time})^{-1}$ and $\beta$ has units inverse time.
Some other conventions use an infection term divided by population size;
we consistently follow the supplied equations without that divisor.

Take $S(0)=S_0>0$, $I(0)=I_0>0$ and $R(0)=0$.
Adding equations gives $S+I+R=N=S_0+I_0$, a conserved total.
Thus $R=N-S-I$ and the physical $S$--$I$ phase region is the triangle
\[
S\geq0,\qquad I\geq0,\qquad S+I\leq N.
\]
The full line $I=0$ consists of equilibria. The $S$-nullclines are $S=0$
and $I=0$; the $I$-nullclines are $I=0$ and $S=\rho$, where
\[
\rho=\frac\beta\alpha.
\]
Inside the positive triangle, $\dot S<0$, so every trajectory moves left.
If $S>\rho$, $\dot I>0$; if $S<\rho$, $\dot I<0$.
The dimensionless initial growth threshold is $\alpha S_0/\beta=S_0/\rho$.

\subsection{An exact phase-plane relation}
For $S,I>0$,
\[
\frac{dI}{dS}=-1+\frac\rho S.
\]
Integrating and applying initial conditions gives
\begin{equation}\label{eq:applied-sir-phase}
I(S)=N-S+\rho\log\frac S{S_0}.
\end{equation}
This relates the two populations without explicitly solving for time.
% !TeX root = ../applied-mathematics1-companion.tex
% Dimensionless illustration N=1, S0=.99, I0=.01, rho=.3; not observed data.
\begin{center}
\begin{tikzpicture}[x=7.4cm,y=5cm,>=Stealth]
\draw[->] (0,0)--(1.06,0) node[right] {$S/N$};
\draw[->] (0,0)--(0,1.06) node[above] {$I/N$};
\draw[muted,dotted] (0,1)--(1,0);
\draw[muted,dashed] (.3,0)--(.3,.7);
\draw[very thick,ink,domain=0.0404097352:.99,samples=100]
 plot (\x,{1-\x+.3*ln(\x/.99)});
\draw[->,thick,ink] (.72,{1-.72+.3*ln(.72/.99)})
 --(.66,{1-.66+.3*ln(.66/.99)});
\draw[->,thick,ink] (.15,{1-.15+.3*ln(.15/.99)})
 --(.12,{1-.12+.3*ln(.12/.99)});
\fill (.99,.01) circle (1.5pt) node[above left,yshift=2pt] {initial};
\fill (.3,{.7+.3*ln(.3/.99)}) circle (1.5pt) node[above right] {peak};
\node[below] at (.3,0) {$\rho/N$};
\node[above right,muted] at (.68,.32) {$S+I=N$};
\node[below left] at (0,0) {$0$};
\end{tikzpicture}

{\small Illustrative SIR trajectory; arrows move towards decreasing $S$.}
\end{center}


\subsection{When does the infection peak?}
If $S_0>\rho$, $I$ initially rises. Its subsequent peak occurs when
$S=\rho$, giving
\[
I_{\max}=N-\rho+\rho\log\frac\rho{S_0}.
\]
If $S_0\leq\rho$, the trajectory starts on or to the left of this threshold,
so its maximum for $t\geq0$ is $I_0$, not the value of an unreachable
turning point obtained by extending the phase curve backwards.
The peak formula therefore needs the outbreak condition $S_0>\rho$.

\subsection{Removed population and final size}
Taking the ratio of the $S$ and $R$ equations gives
\[
\frac{dS}{dR}=-\frac S\rho
\quad\Longrightarrow\quad S(R)=S_0e^{-R/\rho}.
\]
Since $0\leq R\leq N$,
$S(t)\geq S_0e^{-N/\rho}>0$. Susceptibles are not exhausted completely;
infections stop growing when their reproduction rate falls below removals.
Substitution into $\dot R=\beta I$ gives the exact scalar equation
\begin{equation}\label{eq:applied-sir-R}
\dot R=\beta\bigl[N-R-S_0e^{-R/\rho}\bigr],\qquad R(0)=0.
\end{equation}
It is suitable for numerical integration but does not generally have a
simple elementary explicit solution.

\textit{Supplementary final-size relation.} Since $S$ decreases and $R$
increases while both remain bounded, they have limits $S_\infty,R_\infty$.
Then $I=N-S-R$ has a limit too. That limit must be zero, since
$\dot R=\beta I$ and a positive limiting $I$ would force unbounded $R$.
Therefore
\[
S_\infty=S_0e^{-(N-S_\infty)/\rho},\qquad R_\infty=N-S_\infty.
\]
This implicit equation determines the physical final state reached from
the given initial conditions; it is not an instruction to choose an
arbitrary algebraic root.

\section{An approximate explicit SIR solution}
The sources obtain a solvable approximation by assuming $R/\rho\ll1$.
Expanding the exponential in \eqref{eq:applied-sir-R} gives
\[
e^{-R/\rho}=1-\frac R\rho+\frac{R^2}{2\rho^2}
+O\bigl((R/\rho)^3\bigr).
\]
With $I_0=N-S_0$, define $q=S_0/\rho-1$ and $a=S_0/(2\rho^2)$.
The quadratic approximation is
\begin{equation}\label{eq:applied-sir-approx}
\dot R_{\mathrm{app}}=\beta(I_0+qR_{\mathrm{app}}-aR_{\mathrm{app}}^2),
\qquad R_{\mathrm{app}}(0)=0.
\end{equation}
Using a separate subscript emphasises that this is a different differential
equation, approximating the exact model only where the expansion is accurate.

\subsection{Solving the quadratic equation}
The source marks the detailed integration non-examinable. Retain it as an
optional worked derivation. Put
\[
\kappa=\sqrt{q^2+4aI_0},\qquad
b=\frac q{2a},\qquad A=\frac\kappa{2a}.
\]
Since $I_0>0$, $\kappa>|q|$ and hence $|q/\kappa|<1$.
Completing the square gives
\[
I_0+qR-aR^2=a\bigl[A^2-(R-b)^2\bigr].
\]
Let $z=(R-b)/A$. Then
\[
\dot z=\frac{\beta\kappa}{2}(1-z^2),\qquad
\operatorname{artanh}z=\frac{\beta\kappa t}{2}+C.
\]
At $t=0$, $z=-q/\kappa$, so set
$\phi=\operatorname{artanh}(q/\kappa)$ to obtain
\begin{equation}\label{eq:applied-sir-tanh}
R_{\mathrm{app}}(t)=\frac{\rho^2}{S_0}
\left[q+\kappa\tanh\left(\frac{\beta\kappa t}{2}-\phi\right)\right].
\end{equation}
Here $\operatorname{artanh}$ is the inverse hyperbolic tangent, not the
reciprocal of $\tanh$. Differentiating the approximate solution gives
\[
\dot R_{\mathrm{app}}(t)=\frac{\beta\kappa^2\rho^2}{2S_0}
\operatorname{sech}^2\left(\frac{\beta\kappa t}{2}-\phi\right).
\]
The initial condition is a useful check: $\tanh(-\phi)=-q/\kappa$, so
$R_{\mathrm{app}}(0)=0$.

\begin{remark}[Keep the approximation consistent]
For the exact model, $S=S_0e^{-R/\rho}$ and $I=N-R-S$.
Inserting $R_{\mathrm{app}}$ into these expressions gives approximate
reconstructions, not an exact SIR solution. Alternatively
$I_{\mathrm{app}}=\dot R_{\mathrm{app}}/\beta$ follows from the quadratic
truncation. These two reconstructions agree only to the retained order:
they should not be described as identical exact functions.
A large numerical $\rho$ is not sufficient by itself; the dimensionless
ratio $R/\rho$ must remain small over the time range of interest.
\end{remark}

\subsection{What the source's historical illustrations show}
For the supplied boarding-school illustration, the stated inputs are
$N=760$, $I_0=1$, $S_0=759$ and a peak near $300$ infections.
Using the exact phase-plane peak formula gives an equation for $\rho$:
\[
760-\rho+\rho\log(\rho/759)=300.
\]
On $0<\rho<759$, numerical solution gives roughly $\rho\approx195.0$.
The source rounds this to about $200$. The peak height determines this
parameter ratio, but not the time scale: one needs timing information to
identify $\beta$, then $\alpha=\beta/\rho$.

For the source's rounded choice $\rho=200$, $q\approx2.795$,
$\kappa\approx2.802$, and $q/\kappa\approx0.998$.
Thus $\phi$ is around $3.4$, not $\operatorname{artanh}(0.014)$ as printed.
Also the total number removed can be much larger than $200$, so the
small-$R/\rho$ approximation cannot be assumed valid over the whole outbreak.
These are model illustrations from the handout; the underlying data are
not supplied as a table, and no data fit has been reproduced here.

The Bombay plague illustration similarly compares a curve proportional
to $\dot R$, rather than cumulative $R$, with deaths per unit time.
Interpreting deaths as a fixed fraction of removals requires an extra
observation assumption; the compartment equations alone do not provide
that conversion. The source's plotted observations and fitted parameters
are not reconstructed as new evidence.

\section{A second-order equation as a planar system}
For
\[
\ddot x+3\dot x+2\sin x=0,
\]
set $y=\dot x$. Then
\[
\dot x=y,\qquad\dot y=-3y-2\sin x.
\]
The nullclines are $y=0$ and $y=-(2/3)\sin x$, and the equilibria are
$(n\pi,0)$ for integers $n$. At such a point,
\[
J=\begin{pmatrix}0&1\\-2\cos(n\pi)&-3\end{pmatrix},\qquad
\lambda^2+3\lambda+2\cos(n\pi)=0.
\]
For even $n$, the eigenvalues are $-1,-2$, giving local attraction.
For odd $n$, they are $(-3\pm\sqrt{17})/2$, with opposite signs,
giving an unstable saddle.
This conversion lets the same equilibrium and Jacobian methods analyse
second-order motion, connecting this chapter with mechanics.

\studyguide{For a continuous model: find zero rates, not intersections
with $y=x$. Use derivative signs in one dimension and Jacobian eigenvalue
real parts in two dimensions. If the linear test has zero real parts,
look for nonlinear information such as a first integral. State any domain
or small-parameter assumptions before using an explicit approximation.}

% !TeX root = ../applied-mathematics1-companion.tex
\chapter{Mechanics: Simple Forces}\label{chap:applied-mechanics}

\section{Newton's laws and inertial frames}
The models in this chapter treat bodies as point particles of constant mass.
Position, velocity, acceleration and force are vectors; in one-dimensional
motion we use signed scalar components along a chosen positive axis.

\begin{definition}[Inertial frame]
An inertial frame is a reference frame in which a force-free particle moves
with constant velocity. Newton's laws below are used in such a frame.
\end{definition}
\begin{enumerate}
\item \textbf{First law.} An isolated particle moves in a straight line at
constant velocity. Zero velocity is just one possibility.
\item \textbf{Second law.} For fixed mass $m$, the net force is
$\mathbf F_{\mathrm{net}}=m\mathbf a=m\ddot{\mathbf r}$.
\item \textbf{Third law.} The force on particle 1 due to particle 2 equals
minus the force on particle 2 due to particle 1:
$\mathbf F_{1\leftarrow2}=-\mathbf F_{2\leftarrow1}$.
The strong form used for central interactions also places them along the
line joining the particles.
\end{enumerate}
Third-law force pairs act on different bodies: they do not cancel on a
free-body diagram of just one body. Multiple forces acting on the same body
must be summed before applying the second law.
In an accelerating frame, extra inertial-force terms are needed; one cannot
simply reuse the same net-force equation unchanged.

\section{Balancing forces}
Near Earth's surface, weight has magnitude $mg$, directed downwards, with
$g\approx9.8\ \mathrm{m\,s^{-2}}$. Mass is measured in kilograms, force in
newtons, and $1\ \mathrm N=1\ \mathrm{kg\,m\,s^{-2}}$.
Weight depends on the gravitational field; mass does not disappear when
an object is weightless.

A taut light string exerts a tension along itself, pulling on the body to
which it is attached. A smooth contact surface exerts a normal reaction
perpendicular to the surface. For a light string over an ideal smooth
massless pulley, the tension magnitude is the same on both sides.
A string cannot push; if a calculation requires negative tension, the taut
string model must be reconsidered.

Static equilibrium requires zero acceleration and zero initial velocity.
Zero net force alone gives constant velocity, which need not be rest.
To balance forces on a particle, resolve the vector equation
$\sum\mathbf F=\mathbf0$ into convenient perpendicular directions.

\studyguide{Draw a separate free-body diagram for each body or junction.
Choose axes and mark the angles from those axes. Resolve every force,
then apply $\sum F=ma$ along each axis. Do not add a body's own third-law
reaction force to its own diagram.}

\begin{example}[Three supporting cables: source Example 3.1]
A mass $m$ hangs from a vertical cable, attached at a light junction to two
inclined cables making angles $\theta_1,\theta_2$ above the horizontal.
Let the inclined tensions be $T_1,T_2$ and the vertical one $T_3$.
% !TeX root = ../applied-mathematics1-companion.tex
\begin{center}
\begin{tikzpicture}[>=Stealth,scale=.9]
\draw[muted,dashed] (-2.4,0)--(2.6,0);
\draw[->,very thick,ink] (0,0)--(-2,1.5) node[above] {$T_1$};
\draw[->,very thick,ink] (0,0)--(2.2,1) node[above] {$T_2$};
\draw[->,very thick,accent] (0,0)--(0,-1.6) node[right] {$T_3=mg$};
\draw (-.7,0) arc[start angle=180,end angle=143.13,radius=.7];
\draw (.8,0) arc[start angle=0,end angle=24.44,radius=.8];
\node at (-1,.29) {$\theta_1$};
\node at (1.1,.22) {$\theta_2$};
\fill (0,0) circle (2pt);
\end{tikzpicture}

{\small Free-body diagram of the light junction; angles are from horizontal.}
\end{center}

For the mass, equilibrium gives $T_3=mg$.
For the junction, resolve horizontally and vertically:
\[
T_1\cos\theta_1=T_2\cos\theta_2,\qquad
T_1\sin\theta_1+T_2\sin\theta_2=mg.
\]
Eliminating one tension and using the sine addition identity gives
\[
T_1=\frac{mg\cos\theta_2}{\sin(\theta_1+\theta_2)},\qquad
T_2=\frac{mg\cos\theta_1}{\sin(\theta_1+\theta_2)},\qquad T_3=mg.
\]
For equal angles $\theta$, both inclined tensions equal $mg/(2\sin\theta)$.
Nearly horizontal cables need very large tension, consistent with the
small vertical components of their forces.
\end{example}

\begin{example}[Blocks on smooth inclines: source Example 3.2]
Two masses $m_1=1\ \mathrm{kg}$ and $m_2=\sqrt2\ \mathrm{kg}$ are joined
by a taut light string over an ideal pulley. The smooth inclines have
angles $\pi/4$ and $\theta$ above horizontal.
For equilibrium, resolve along each plane:
\[
T=m_1g\sin(\pi/4)=\frac g{\sqrt2},\qquad
T=m_2g\sin\theta=\sqrt2g\sin\theta.
\]
Thus $\sin\theta=1/2$; with $0<\theta<\pi/2$, $\theta=\pi/6$.
Perpendicular resolution gives
\[
R_1=m_1g\cos(\pi/4)=\frac g{\sqrt2},\qquad
R_2=m_2g\cos(\pi/6)=g\sqrt{\frac32}.
\]
With masses substituted numerically in kilograms, these force values are
in newtons. The reaction is not generally $mg$: only the perpendicular
component of weight is balanced by it.
\end{example}

\section{Hooke's law and connected springs}
A spring of natural length $\ell$ and stiffness $k>0$ exerts a restoring
force proportional to signed extension $e$:
\[
\text{spring tension}=ke,\qquad\text{force on a displaced end opposes }e.
\]
Stiffness has units $\mathrm{N/m}$. Positive extension means stretching;
negative extension means compression, when a spring pushes outwards.
Hooke's law is an approximation within the spring's elastic range.
An elastic string differs: it supplies tension when extended but goes
slack when shorter than its natural length.

\begin{example}[Three vertical springs: source Example 3.3]
Three identical light springs of natural length $\ell$ and stiffness $k$
connect a ceiling to a floor $3\ell$ below it. A mass $m$ is attached at
each of the two intermediate junctions.
Measure distance downwards from the ceiling. Let the mass positions be
$\ell+x_1$ and $2\ell+x_2$. The signed spring extensions are
\[
e_1=x_1,\qquad e_2=x_2-x_1,\qquad e_3=-x_2.
\]
For the upper mass, downward forces minus upward forces give
\[
mg+k(x_2-x_1)-kx_1=0.
\]
For the lower mass, the middle spring force is upward with signed magnitude
$k(x_2-x_1)$, and the bottom spring contributes downward signed force
$-kx_2$. Thus
\[
mg-k(x_2-x_1)-kx_2=0.
\]
The two equations become
\[
2x_1-x_2=\frac{mg}{k},\qquad2x_2-x_1=\frac{mg}{k}.
\]
Subtracting gives $x_1=x_2$, and hence both equal $mg/k$.
The mass positions below the ceiling are
\[
\ell+\frac{mg}{k},\qquad2\ell+\frac{mg}{k}.
\]
The spring lengths are $\ell+mg/k$, $\ell$, and $\ell-mg/k$.
The bottom spring is compressed and pushes upwards on the lower mass;
the middle spring has zero extension. The ideal geometry requires
$mg/k<\ell$, and physical validity also requires staying in the elastic range.
\end{example}

\section{Equations of one-dimensional motion}
Choose a signed coordinate $x(t)$ and write
\[
v=\dot x,\qquad a=\ddot x=\frac{dv}{dt},\qquad
m\ddot x=F(t,x,\dot x).
\]
Where velocity can be considered locally as a function of position, the
chain rule gives a third useful acceleration form:
\[
a=\frac{dv}{dt}=\frac{dv}{dx}\frac{dx}{dt}=v\frac{dv}{dx}.
\]
Use $dv/dt$ for a velocity--time relation and $v\,dv/dx$ for a
velocity--position relation. Turning points may require treating separate
branches of $v(x)$; one cannot divide by $v$ indiscriminately at $v=0$.

\section{Constant force and time-dependent force}
\subsection{Constant acceleration}
For $\dot v=a$ constant, with $x(0)=x_0$ and $v(0)=u$, integrate twice:
\[
v(t)=u+at,\qquad x(t)=x_0+ut+\frac12at^2.
\]
Alternatively, integrating $v\,dv/dx=a$ gives
\[
v^2=u^2+2a(x-x_0).
\]
These formulae require constant acceleration and consistent signed directions.

\begin{example}[A sinusoidal force: source Example 3.4]
Suppose $F(t)=mA\sin(\omega t)$, with $\omega>0$, $x(0)=0$ and
$v(0)=v_0$. Here $A$ is an acceleration amplitude, not a force.
The equation $\ddot x=A\sin(\omega t)$ integrates to
\begin{align*}
v(t)&=v_0+\frac A\omega\bigl(1-\cos(\omega t)\bigr),\\
x(t)&=\left(v_0+\frac A\omega\right)t
-\frac A{\omega^2}\sin(\omega t).
\end{align*}
The sine term stays bounded, while the first term is linear in $t$.
Consequently $x(t)\to+\infty$ exactly when $v_0>-A/\omega$;
it tends to $-\infty$ for $v_0<-A/\omega$, and remains bounded at equality.
Divergence of position to $+\infty$ does not necessarily mean the velocity
is positive at every sufficiently late instant: bounded oscillations can
remain on top of the drift.
\end{example}

\section{Quadratic resistance: direction matters}
A drag force must oppose velocity. With constant $k>0$, quadratic drag per
unit mass is represented in one dimension by $-kv|v|$.
Choose upward displacement $y$ and upward signed velocity $v=\dot y$.
With gravity, a single signed equation is
\[
\dot v=-g-kv|v|.
\]
During ascent ($v>0$), it becomes $\dot v=-g-kv^2$.
During descent ($v<0$), it becomes $\dot v=-g+kv^2$.
Using $-kv^2$ for both upward and downward velocities would make drag point
downward in both cases, which is physically wrong.
Here $k$ has units inverse length.

\begin{example}[Maximum height: source Example 3.5]
Launch from $y=0$ with $v(0)=u>0$. At the highest point $y=h$, $v=0$.
For the ascending part use $v\,dv/dy=-g-kv^2$, so
\begin{align*}
h&=-\int_u^0\frac{v\,dv}{g+kv^2}\\
&=\frac1{2k}\log\left(\frac{g+ku^2}{g}\right)
=\frac1{2k}\log\left(1+\frac{ku^2}{g}\right).
\end{align*}
As $k\to0$, $\log(1+z)\sim z$ recovers the drag-free height $u^2/(2g)$.
Only ascent was used; the drag sign changes after the turning point.
\end{example}

\begin{example}[Downward terminal speed: source Example 3.6]
Now choose displacement positive \emph{downwards} and let $v(0)=v_0\geq0$.
Gravity is positive and drag negative, so
\[
\dot v=g-kv^2=k(V^2-v^2),\qquad V=\sqrt{\frac gk}.
\]
The equilibrium $v=V$ is the terminal speed. If $v_0=V$, the speed stays
constant. Otherwise, separation gives
\[
\int_{v_0}^{v(t)}\frac{dw}{V^2-w^2}=kt,
\qquad
\frac1{V^2-w^2}=\frac1{2V}\left(\frac1{V-w}+\frac1{V+w}\right).
\]
Hence
\[
\log\left|\frac{V+v}{V-v}\frac{V-v_0}{V+v_0}\right|=2Vkt.
\]
A solution cannot cross the equilibrium $V$ by uniqueness, so $V-v$ keeps
the sign of $V-v_0$. Rearranging gives the expression valid for all
$v_0\geq0$:
\begin{equation}\label{eq:applied-quadratic-drag}
v(t)=V\frac{(V+v_0)e^{2Vkt}-(V-v_0)}
              {(V+v_0)e^{2Vkt}+(V-v_0)}.
\end{equation}
It also includes $v_0=V$ by direct substitution. For $v_0<V$ the speed
increases towards $V$; for $v_0>V$ it decreases towards $V$.
Equivalently, when $v_0<V$,
\[
v(t)=V\tanh\left(Vkt+\operatorname{artanh}(v_0/V)\right).
\]
For $v_0>V$, the corresponding hyperbolic form uses $\coth$, not a real
$\operatorname{artanh}(v_0/V)$. The exponential fraction avoids this case split.
\end{example}

\section{Linear resistance and terminal velocity}
For constant applied force per unit mass $a_0$ and linear drag coefficient
$\lambda>0$, the signed equation is
\[
\dot v=a_0-\lambda v.
\]
The drag term $-\lambda v$ automatically reverses with the velocity.
The terminal velocity is $v_\infty=a_0/\lambda$, which may be positive,
negative or zero. The coefficient $\lambda$ has units inverse time.
An integrating factor gives, for every real $v_0$,
\begin{align*}
v(t)&=v_\infty+(v_0-v_\infty)e^{-\lambda t},\\
x(t)&=x_0+v_\infty t+
\frac{v_0-v_\infty}{\lambda}(1-e^{-\lambda t}).
\end{align*}
The difference from terminal velocity preserves its sign and decays
exponentially. Starting exactly at terminal velocity produces uniform motion.

\begin{example}[A bullet in a block: source Example 3.8]
A particle enters a resistive block at $x=0$ with speed $u>0$. Resistance
per unit mass is $kv$, with $k>0$, and gravity is neglected.
Then $\dot v=-kv$ and
\[
v(t)=ue^{-kt},\qquad x(t)=\frac uk(1-e^{-kt}).
\]
The limiting penetration is $u/k$. It is approached as $t\to\infty$;
in this ideal model, the particle does not stop exactly in finite time.
A block of width at least $u/k$ prevents finite-time escape through the
far side. To reach half the limiting distance,
\[
\frac u{2k}=\frac uk(1-e^{-kT})
\quad\Longrightarrow\quad T=\frac{\log2}{k}.
\]

Now two particles enter simultaneously from opposite ends, initially
separated by $L>0$, with inward speeds $u_1,u_2>0$.
Their distances into the block at time $t$ are
\[
s_1(t)=\frac{u_1}{k}(1-e^{-kt}),\qquad
s_2(t)=\frac{u_2}{k}(1-e^{-kt}).
\]
Finite-time collision requires $s_1+s_2=L$. Since $1-e^{-kt}<1$ for every
finite $t\geq0$, the necessary and sufficient condition is
\[
\boxed{u_1+u_2>kL}.
\]
When $u_1+u_2=kL$, the separation tends to zero but is positive at every
finite time; that is a limiting meeting, not a finite-time collision.
This corrects the non-strict inequality in both supplied versions.
For a finite collision,
\[
T=\frac1k\log\frac{u_1+u_2}{u_1+u_2-kL},\qquad
s_1(T)=\frac{u_1L}{u_1+u_2}.
\]
The second expression gives the collision position measured from the
first particle's entry end. As $u_1+u_2$ decreases to $kL$ from above,
the collision time diverges, confirming the strict threshold.
\end{example}

\studyguide{For motion: choose the positive direction first; draw forces;
write the signed net force; choose $dv/dt$ or $v\,dv/dx$ according to the
relation wanted; include initial conditions and equilibrium velocities;
check the result's units, signs and limiting behaviour.}

\end{document}
